\documentclass[CJK,13pt]{beamer}
\usepackage{CJKutf8}
\usepackage{beamerthemesplit}
\usetheme{Malmoe}
\useoutertheme[footline=authortitle]{miniframes}
\usepackage{amsmath}
\usepackage{amssymb}
\usepackage{graphicx}
\usepackage{eufrak}
\usepackage{color}
\usepackage{slashed}
\usepackage{simplewick}
\usepackage{tikz}
\usepackage{tcolorbox}
\usepackage{ulem}
\graphicspath{{../figures/}}
%%figures
\def\lfig#1#2{\includegraphics[width=#1 in]{#2}}
\def\tfig#1#2{\includegraphics[height=#1 in]{#2}}
\def\addfig#1#2{\begin{center}\includegraphics[width=#1 in]{#2}\end{center}}
\def\question{\includegraphics[width=0.3in]{why.jpg}\,}
\def\answer{\includegraphics[width=0.3in]{baozou_haha.png}\,}
\def\wulian{\includegraphics[width=0.18in]{emoji_wulian.jpg}}
\def\bigwulian{\includegraphics[width=0.35in]{emoji_wulian.jpg}}
\def\bye{\includegraphics[width=0.18in]{emoji_bye.jpg}}
\def\bigbye{\includegraphics[width=0.35in]{emoji_bye.jpg}}
\def\huaixiao{\includegraphics[width=0.18in]{emoji_huaixiao.jpg}}
\def\bighuaixiao{\includegraphics[width=0.35in]{emoji_huaixiao.jpg}}
\def\jianxiao{\includegraphics[width=0.18in]{emoji_jianxiao.jpg}}
\def\bigjianxiao{\includegraphics[width=0.35in]{emoji_jianxiao.jpg}}
\def\haoqi{\includegraphics[width=0.18in]{emoji_haoqi.jpg}}
%% colors
\def\blacktext#1{{\color{black}#1}}
\def\bluetext#1{{\color{blue}#1}}
\def\redtext#1{{\color{red}#1}}
\def\darkbluetext#1{{\color[rgb]{0,0.2,0.6}#1}}
\def\skybluetext#1{{\color[rgb]{0.2,0.7,1.}#1}}
\def\cyantext#1{{\color[rgb]{0.,0.5,0.5}#1}}
\def\greentext#1{{\color[rgb]{0,0.7,0.1}#1}}
\def\darkgray{\color[rgb]{0.2,0.2,0.2}}
\def\lightgray{\color[rgb]{0.6,0.6,0.6}}
\def\gray{\color[rgb]{0.4,0.4,0.4}}
\def\blue{\color{blue}}
\def\red{\color{red}}
\def\orange{\color[rgb]{1.,0.8,0.}}
\def\green{\color{green}}
\def\darkgreen{\color[rgb]{0,0.4,0.1}}
\def\darkblue{\color[rgb]{0,0.2,0.6}}
\def\skyblue{\color[rgb]{0.2,0.7,1.}}
%%control
\def\diag{\mathrm{diag}\,}
\def\heaviside{\,\mathrm{h}\,}
\def\bral{\left(\begin{array}{l}}
\def\brar{\end{array}\right)}
\def\brall{\left(\begin{array}{ll}}
\def\brarr{\end{array}\right)}
\def\bralll{\left(\begin{array}{lll}}
\def\brarrr{\end{array}\right)}
\def\branchl{\left\{\begin{array}{l}}
\def\branchr{\end{array}\right.}
\def\branchll{\left\{\begin{array}{ll}}
\def\branchrr{\end{array}\right.}
\def\branchlll{\left\{\begin{array}{lll}}
\def\branchrrr{\end{array}\right.}
\def\sfgamma#1{\,\Gamma\left( #1 \right)\,}
\def\be{\begin{equation}}
\def\ee{\nonumber\end{equation}}
\def\bea{\begin{eqnarray}}
\def\eea{\nonumber\end{eqnarray}}
\def\bch{\begin{CJK}{UTF8}{gbsn}}
\def\ech{\end{CJK}}
\def\bitem{\begin{itemize}}
\def\eitem{\end{itemize}}
\def\bcenter{\begin{center}}
\def\ecenter{\end{center}}
\def\bex{\begin{minipage}{0.2\textwidth}\includegraphics[width=0.6in]{jugelizi.png}\end{minipage}\begin{minipage}{0.76\textwidth}}
\def\eex{\end{minipage}}
\def\chtitle#1{\frametitle{\bch#1\ech}}
\def\bmat#1{\left(\begin{array}{#1}}
\def\emat{\end{array}\right)}
\def\bcase#1{\left\{\begin{array}{#1}}
\def\ecase{\end{array}\right.}
\def\bmini#1{\begin{minipage}{#1\textwidth}}
\def\emini{\end{minipage}}
\def\tbox#1{\begin{tcolorbox}#1\end{tcolorbox}}
\def\pfrac#1#2#3{\left(\frac{\partial #1}{\partial #2}\right)_{#3}}
\def\res#1#2{\,\mathrm{res}\,#1\left(#2\right)\,}
\def\newt#1#2{\left(\begin{array}{c}#1\\ #2\end{array}\right)}
\def\reof#1{\,\mathrm{Re}{\left(#1\right)\,}}
\def\imof#1{\,\mathrm{Im}{\left(#1\right)\,}}
\def\Arg#1{\,\mathrm{Arg}\,#1\,}
%%symbols
\def\intfull{\int_{-\infty}^\infty}
\def\inthalf{\int_0^\infty}
\def\bropt{\,(\ \ \ )}
\def\sech{\mathrm{sech}\,}
\def\csch{\mathrm{csch}\,}
\def\asinh{\mathrm{asinh}\,}
\def\acosh{\mathrm{acosh}\,}
\def\sone{$\star$}
\def\stwo{$\star\star$}
\def\sthree{$\star\star\star$}
\def\sfour{$\star\star\star\star$}
\def\sfive{$\star\star\star\star\star$}
\def\rint{{\int_\leftrightarrow}}
\def\roint{{\oint_\leftrightarrow}}
\def\stdHf{{\textit{\r H}_f}}
\def\deltaH{{\Delta \textit{\r H}}}
\def\ii{{\dot{\imath}}}
\def\skipline{{\vskip0.1in}}
\def\skiplines{{\vskip0.2in}}
\def\lagr{{\mathcal{L}}}
\def\hamil{{\mathcal{H}}}
\def\vecT{{\mathbf{T}}}
\def\vecN{{\mathbf{N}}}
\def\vecB{{\mathbf{B}}}
\def\vecv{{\mathbf{v}}}
\def\vecr{{\mathbf{r}}}
\def\vecf{{\mathbf{f}}}
\def\vecg{{\mathbf{g}}}
\def\vecupsilon{{\mathbf{\upsilon}}}
\def\vecu{{\mathbf{u}}}
\def\vecj{{\mathbf{j}}}
\def\vecx{{\mathbf{x}}}
\def\vecy{{\mathbf{y}}}
\def\vecz{{\mathbf{z}}}
\def\veck{{\mathbf{k}}}
\def\vecp{{\mathbf{p}}}
\def\vecn{{\mathbf{n}}}
\def\vecA{{\mathbf{A}}}
\def\vecP{{\mathbf{P}}}
\def\vecsigma{{\mathbf{\sigma}}}
\def\hatJn{{\hat{J_\vecn}}}
\def\hatJx{{\hat{J_x}}}
\def\hatJy{{\hat{J_y}}}
\def\hatJz{{\hat{J_z}}}
\def\hatj#1{\hat{J_{#1}}}
\def\hatphi{{\hat{\phi}}}
\def\hatq{{\hat{q}}}
\def\hatpi{{\hat{\pi}}}
\def\vel{\upsilon}
\def\Dint{{\mathcal{D}}}
\def\adag{{\hat{a}^\dagger}}
\def\bdag{{\hat{b}^\dagger}}
\def\cdag{{\hat{c}^\dagger}}
\def\ddag{{\hat{d}^\dagger}}
\def\hata{{\hat{a}}}
\def\hatb{{\hat{b}}}
\def\hatc{{\hat{c}}}
\def\hatd{{\hat{d}}}
\def\hatD{{\,\hat{D}}}
\def\hatN{{\hat{N}}}
\def\hatH{{\hat{H}}}
\def\hatp{{\hat{p}}}
\def\Fup{{F^{\mu\nu}}}
\def\Fdown{{F_{\mu\nu}}}
\def\newl{\nonumber \\}
\def\vece{\mathrm{e}}
\def\calM{{\mathcal{M}}}
\def\calT{{\mathcal{T}}}
\def\calR{{\mathcal{R}}}
\def\barpsi{\bar{\psi}}
\def\baru{\bar{u}}
\def\barv{\bar{\upsilon}}
\def\qeq{\stackrel{?}{=}}
\def\ftf{\stackrel{\mathcal{FT}}{\Longrightarrow}}
\def\ftb{\stackrel{\mathcal{FT}}{\Longleftarrow}}
\def\ftfb{\stackrel{\mathcal{FT}}{\Longleftrightarrow}}
\def\ltf{\stackrel{\mathcal{LT}}{\Longrightarrow}}
\def\ltb{\stackrel{\mathcal{LT}}{\Longleftarrow}}
\def\ltfb{\stackrel{\mathcal{LT}}{\Longleftrightarrow}}
\def\torder#1{\mathcal{T}\left(#1\right)}
\def\rorder#1{\mathcal{R}\left(#1\right)}
\def\contr#1#2{\contraction{}{#1}{}{#2}#1#2}
\def\trof#1{\mathrm{Tr}\left(#1\right)}
\def\trace{\mathrm{Tr}}
\def\comm#1{\ \ \ \left(\mathrm{used}\ #1\right)}
\def\tcomm#1{\ \ \ (\text{#1})}
\def\slp{\slashed{p}}
\def\slk{\slashed{k}}
\def\calp{{\mathfrak{p}}}
\def\veccalp{\mathbf{\mathfrak{p}}}
\def\Tthree{T_{\tiny \textcircled{3}}}
\def\pthree{p_{\tiny \textcircled{3}}}
\def\dbar{{\,\mathchar'26\mkern-12mu d}}
\def\erf{\mathrm{erf}}
\def\const{\mathrm{const.}}
\def\pheat{\pfrac p{\ln T}V}
\def\vheat{\pfrac V{\ln T}p}
%%units
\def\fdeg{{^\circ \mathrm{F}}}
\def\cdeg{^\circ \mathrm{C}}
\def\atm{\,\mathrm{atm}}
\def\angstrom{\,\text{\AA}}
\def\SIL{\,\mathrm{L}}
\def\SIkm{\,\mathrm{km}}
\def\SIyr{\,\mathrm{yr}}
\def\SIGyr{\,\mathrm{Gyr}}
\def\SIV{\,\mathrm{V}}
\def\SImV{\,\mathrm{mV}}
\def\SIeV{\,\mathrm{eV}}
\def\SIkeV{\,\mathrm{keV}}
\def\SIMeV{\,\mathrm{MeV}}
\def\SIGeV{\,\mathrm{GeV}}
\def\SIcal{\,\mathrm{cal}}
\def\SIkcal{\,\mathrm{kcal}}
\def\SImol{\,\mathrm{mol}}
\def\SIN{\,\mathrm{N}}
\def\SIHz{\,\mathrm{Hz}}
\def\SIm{\,\mathrm{m}}
\def\SIcm{\,\mathrm{cm}}
\def\SIfm{\,\mathrm{fm}}
\def\SImm{\,\mathrm{mm}}
\def\SInm{\,\mathrm{nm}}
\def\SImum{\,\mathrm{\mu m}}
\def\SIJ{\,\mathrm{J}}
\def\SIW{\,\mathrm{W}}
\def\SIkJ{\,\mathrm{kJ}}
\def\SIs{\,\mathrm{s}}
\def\SIkg{\,\mathrm{kg}}
\def\SIg{\,\mathrm{g}}
\def\SIK{\,\mathrm{K}}
\def\SImmHg{\,\mathrm{mmHg}}
\def\SIPa{\,\mathrm{Pa}}
%page
\def\secpage#1#2{\begin{frame}\bcenter{\bf \Huge #1} \ecenter \skipline \tbox{\bcenter #2 \ecenter}\end{frame}}
\def\append#1#2{\secpage{附录#1}{#2}}
\def\thinka#1{\begin{frame}\frametitle{思考题}\bcenter\lfig{0.4}{think0.jpg}  \skipline #1 \ecenter\end{frame}}
\def\thinkb#1{\begin{frame}\frametitle{思考题}\bcenter\lfig{0.5}{think1.jpg}  \skipline #1 \ecenter\end{frame}}
\def\thinkc#1{\begin{frame}\frametitle{思考题}\bcenter\lfig{0.5}{think2.jpg}  \skipline #1 \ecenter\end{frame}}
\def\thinkd#1{\begin{frame}\frametitle{思考题}\bcenter\lfig{0.5}{think3.jpg}  \skipline #1 \ecenter\end{frame}}
\def\thinke#1{\begin{frame}\frametitle{思考题}\bcenter\lfig{0.5}{think4.jpg}  \skipline #1 \ecenter\end{frame}}
\def\thinkf#1{\begin{frame}\frametitle{思考题}\bcenter\lfig{0.5}{think5.jpg}  \skipline #1 \ecenter\end{frame}}
\def\schw{ds^2 = \left(1-\frac{2GM}{r}\right)dt^2 - \left(1-\frac{2GM}{r}\right)^{-1} dr^2 - r^2\left(d\theta^2 + \sin^2\theta d\phi^2\right)}

\input{titlepage.tex}
  \date{}
  \begin{document}
  \bch
  \tpage{23}{Gravitational Radiation (II)}

  \def\Tcal{\mathcal{T}}
  \def\Pcal{\mathcal{P}}
  \def\Mcal{\mathcal{M}}
  
  \begin{frame}
    这一讲的主题是：用四极矩近似算引力波

    \addfig{1}{sijiju.jpg}
  \end{frame}


  \secpage{张量的 virial 定理}{$$ \frac{\partial^2}{\partial t^2} \int x^i x^j T^{00}(\vecx, t) d^3\vecx  = 2\int T^{ij}  d^3\vecx$$}

  \begin{frame}

    在Minkowski 时空，对局域、守恒的能量动量张量 $T^{\mu\nu}(\vecx, t)$，定义能量密度的四极矩：
    $$ Q^{ij}(\vecx, t)\equiv \int x^i x^j T^{00}(\vecx, t) d^3\vecx $$
    由局域条件知道三维散度的积分
    $$ \int \partial_k\left(T^{k0}x^ix^j\right) = 0$$
    于是
    {\small
    \be
    \frac{\partial Q^{ij}}{\partial t} = \int \partial_\mu\left(T^{\mu 0}x^ix^j\right)d^3\vecx = \int  T^{\mu 0}\partial_{\mu}(x^ix^j)d^3\vecx= \int \left(T^{i 0}x^j + T^{j0}x^i\right)d^3\vecx 
    \ee}
  \end{frame}


  \begin{frame}

    故技重施，由三维散度积分
    $$\int \partial_k \left(T^{ik}x^j\right) d^3\vecx =0$$
    可以得到{\small
    $$\frac{\partial}{\partial t} \int T^{i0}x^j d^3\vecx = \int \partial_\mu \left(T^{i\mu}x^j\right) d^3\vecx = \int T^{i\mu} \partial_\mu x^j d^3\vecx = \int T^{ij}d^3\vecx $$    
    }
    于是{\blue
      $$\frac{\partial^2 Q^{ij}}{\partial t^2} = 2\int T^{ij}d^3\vecx $$
      这叫做（张量的）virial定理。
    }
  \end{frame}
  

  \secpage{空空形式的引力辐射公式}{$$\frac{dP}{d\Omega} =  \frac{G \omega^2}{\pi } \Mcal_{ij}^*(\omega\vecn)\Mcal_{kl}(\omega\vecn) \Pcal^{ijkl}(\vecn)$$}
  
  \begin{frame}
    为了叙述简明，我们考虑连续谱的情形
    $$\frac{dE}{d\Omega d\omega}\left(\vecn\right) =  \frac{G\omega^2}{4\pi^2}\left(\mathcal{T}^*_{\mu\nu}(k)\mathcal{T}^{\mu\nu}(k)-\frac{1}{2}\lvert \mathcal{T}^\alpha_{\ \alpha}(k)\rvert^2\right)$$    
    守恒条件 $T_{\mu\nu}^{\ \ \ ;\mu} = 0$ 在傅立叶空间表述为
    $$ k^\mu \mathcal{T}_{\mu\nu} = 0$$
    这样四个方程就允许我们{\bf 把带时间指标的分量全都用只带空间指标的分量表示出来}。下面我们一律在傅立叶空间讨论，且省略宗量 $k=(\omega,\omega\vecn)$
  \end{frame}

  \begin{frame}
    \bea
    \mathcal{T}_{0i}  &=& -n^j \mathcal{T}_{ij}, \mathcal{T}^{0i} = n^j\mathcal{T}_{ij} \newl
    \mathcal{T}_{00} &=& \mathcal{T}^{00} = n^in^j \mathcal{T}_{ij} \newl
    \mathcal{T}^{\alpha}_{\ \alpha} &=& \left(n^in^j - \delta^{ij}\right) \mathcal{T}_{ij}
    \eea
    按照惯例这里的拉丁字母 $i, j$ 只对空间指标 $1,2,3$ 求和。
    于是
    \bea
    \lvert \mathcal{T}^\alpha_{\ \alpha}\rvert^2 &=& \left(n^in^j - \delta^{ij}\right) \left(n^kn^l - \delta^{kl}\right)  \mathcal{T}_{ij}^* \mathcal{T}_{kl} \newl
  %  \Tcal_{00}^*  \Tcal^{00} &=& n^in^jn^kn^l \Tcal_{ij}^*\Tcal_{kl} \newl
  %  \Tcal_{0i}^*  \Tcal^{0i}+\Tcal_{i0}^*  \Tcal^{i0} &=& 2n^in^k\delta^{jl} \Tcal_{ij}^*\Tcal_{kl} \newl
  %  \Tcal_{ij}^*\Tcal^{ij} &=& \delta^{ik}\delta^{jl}\Tcal_{ij}^*\Tcal_{kl} \newl
    \Tcal^*_{\mu\nu}\Tcal^{\mu\nu} &=& \left(n^in^k-\delta^{ik}\right)\left(n^jn^l-\delta^{jl}\right)\Tcal_{ij}^*\Tcal_{kl}
    \eea
  \end{frame}

  \begin{frame}
    \frametitle{用空空分量表示的引力波辐射公式}
    {\blue
    \be
    \frac{dE}{d\Omega d\omega}\left(\vecn\right) = \frac{G\omega^2}{4\pi^2} \Tcal_{ij}^*(k)\Tcal_{kl}(k) \Pcal^{ijkl}(\vecn)
    \ee}
    这里的{\blue
    \be
    \Pcal^{ijkl}(\vecn) \equiv \left(n^in^k-\delta^{ik}\right)\left(n^jn^l-\delta^{jl}\right) - \frac{1}{2}\left(n^in^j - \delta^{ij}\right) \left(n^kn^l - \delta^{kl}\right)
    \ee}
    叫做 {\blue 横向无迹投影算符}。

    \skipline
    
    单频的情况也类似，就不重复推导了：
    {\blue
      \be
      \frac{dP}{d\Omega} =  \frac{G \omega^2}{\pi } \Mcal_{ij}^*(\omega\vecn)\Mcal_{kl}(\omega\vecn) \Pcal^{ijkl}(\vecn)
      \ee      
    }
  \end{frame}


    \secpage{四极矩近似：波源尺寸远小于波长}{$$\frac{dP}{d\Omega} \approx  \frac{G \omega^6}{4\pi } Q_{ij}^*Q_{kl} \Pcal^{ijkl}(\vecn)$$}

    \begin{frame}
      \frametitle{四极矩近似的物理场景}
      如果引力波源的运动是非相对论的，那么波源的尺度很可能比引力波的波长小得多（因两者频率相同，而引力波以光速传播）。

      \skipline

      这时就又有了操作的空间——

      \addfig{1}{kaishibiaoyan.jpg}
    \end{frame}


    \begin{frame}
      \frametitle{以单频引力波辐射为例}
      如果波源的尺度远小于 $\frac{1}{\omega}$。那么对波源的振幅 $M^{ij}(\vecx)$ 做傅立叶变换时， $e^{i\veck\cdot \vecx}$ 可以近似当成 $1$。
      $$ \Mcal^{ij}(\omega\vecn)\approx \int  M^{ij}(\vecx) d^3\vecx$$
      对张量 $M^{\mu\nu}(\vecx)e^{-i\omega t}$ 使用virial定理
      $$\int M^{ij}(\vecx) d^3\vecx  = -\frac{1}{2}\omega^2 \int  x^ix^jM^{00}(\vecx) d^3\vecx $$
    \end{frame}

    \begin{frame}
      \frametitle{单频四极矩辐射公式}
      记能量密度的四极矩(quadrupole)
      {\blue      $$ Q^{ij} \equiv \int  x^ix^jM^{00}(\vecx) d^3\vecx $$}
      则 $\Mcal^{ij}(\omega\vecn) \approx \frac{\omega^2}{2}Q^{ij}$。代入到引力波辐射公式里
      {\blue
        \be
        \frac{dP}{d\Omega} \approx  \frac{G \omega^6}{4\pi } Q_{ij}^*Q_{kl} \Pcal^{ijkl}(\vecn)
        \ee
      }
    \end{frame}

    \begin{frame}
      \frametitle{连续谱的四极矩辐射公式}
      对连续谱，推导完全类似。可以定义
      {\blue      $$ Q^{ij}(\omega) \equiv \int  e^{i\omega t}dt\int  x^ix^jT^{00}(t, \vecx) d^3\vecx $$}
      有
      {\blue
      \be
    \frac{dE}{d\Omega d\omega}\left(\vecn\right) \approx \frac{G\omega^6}{16\pi^2} Q_{ij}^*(\omega)Q_{kl}(\omega) \Pcal^{ijkl}(\vecn)
      \ee}
    \end{frame}
    
    \secpage{绕转双星的辐射功率}{$$ P =  \frac{128 GM^2R^4 \omega^6}{5}$$}

    \begin{frame}
      因为只是做个示范，我们考虑最简单的例子：两个质量为 $M$ 的中子星以角频率 $\omega$ 在圆轨道上互相绕转。设它们之间距离为 $2R$，则有
      $$\omega^2 = \frac{GM}{4R^3}$$
      取旋转中心为原点，旋转轴为 $z$ 轴，则近似有{\scriptsize
      $$ T_{00}(t,x, y, z) = M \left[\delta(x-R\cos \omega t)\delta(y-R\sin\omega t)+  \delta(x+R\cos \omega t)\delta(y+R\sin\omega t)\right]\delta (z)$$}
      显然
      $$ Q_{13}=Q_{23}=Q_{33} = 0$$
    \end{frame}

    \begin{frame}
      由于
      $$\int T_{00}(t,x,y,z) x^2 d^3\vecx = MR^2\left[1+\cos(2\omega t)\right] $$
      $$\int T_{00}(t,x,y,z) y^2 d^3\vecx = MR^2\left[1-\cos(2\omega t)\right] $$
      $$\int T_{00}(t,x,y,z) xy d^3\vecx = MR^2\left[\sin(2\omega t)\right] $$            
      所以这是角频率为 $2\omega$ 的单频源，且
      $$ Q_{11} = \frac{MR^2}{2}, Q_{22}=-\frac{MR^2}{2}, Q_{12} = -i \frac{MR^2}{2}$$

    \end{frame}

    \begin{frame}
      代入四极矩辐射公式，得到角频率为 $2\omega$ 的引力波辐射强度为
      $$\frac{dP}{d\Omega} (\vecn) = \frac{2GM^2R^4 \omega^6}{\pi}\left(\sin^4\theta - 8\sin^2\theta + 8\right)$$
      这里的 $\theta$ 是 $\vecn$ 和 $z$ 轴的夹角。

      \skipline
      
      如果对所有方向积分，则得到辐射功率
      $$ P =  \frac{128 GM^2R^4 \omega^6}{5}$$

      这部分的计算是用代码完成的，请参考
      \url{http://zhiqihuang.top/gr/codes/TTproj.py}
    \end{frame}


  \secpage{附录：横向无迹投影算符$\Pcal^{ijkl}$的物理意义}{$$\left(\Tcal^{\rm TT}\right)^{ij} = P^{ijkl}T_{kl}$$}
    \begin{frame}

      把三维欧氏空间的对称二阶张量 $\Tcal_{ij}$ (6个自由度） 分解为2个标量(2个自由度)，一个无源的矢量（2个自由度），和一个横向无迹的二阶张量（2个自由度）：
      $$\Tcal_{ij} = \Phi\delta_{ij} + \left(n_in_j - \frac{1}{3}\delta_{ij}\right)\Psi + (n_iA_j + n_j A_i) + \Tcal^{\rm TT}_{ij} $$
        这里的 $\Phi$, $\Phi$ 是标量； $A_i$ 是无源的三维矢量，满足 $n^iA_i=0$ (现在我在讨论三维欧氏空间，指标在上面和下面都一样)。

        首先两边求迹，可以确定 $\Phi = \frac{\Tcal}{3}$，这里 $\Tcal$ 是 $\Tcal^i_{\ i}$ 的简写。

        然后两边乘以 $n^in^j$，可以得到：
        $$n^in^j\Tcal_{ij} = \frac{\Tcal}{3} + \frac{2}{3}\Psi $$
        即
        $$\Psi = \frac{3}{2}\left(n^in^j\Tcal_{ij} - \frac{\Tcal}{3}\right)$$
  \end{frame}


    \begin{frame}
      然后在
      $$\Tcal_{ij} = \Phi\delta_{ij} + \left(n_in_j - \frac{1}{3}\delta_{ij}\right)\Psi + (n_iA_j + n_j A_i) + \Tcal^{\rm TT}_{ij} $$
      两边乘以 $n^i$ 并代入 $\Phi, \Psi$ 的解，得到
      $$ n^i\Tcal_{ij} = \frac{\Tcal}{3}n_j +  n_j\left(n^kn^l\Tcal_{kl} - \frac{\Tcal}{3}\right) + A_j $$
      即
      $$ A_j =  n^i\Tcal_{ij} -  \left(n^kn^l\Tcal_{kl}\right)n_j $$
      最后，把 $\Phi, \Psi, A_j$ 都代入，得出 $ \left(\Tcal^{\rm TT}\right)^{ij} = P^{ijkl}T_{kl}$.
    \end{frame}


    
    
    
  
    \ech
\end{document} 




  
